\clearpage\SourceLeaf{385}
\pdfbookmark[0]{Sechster Abschnitt: Verschiedene Anwendungen der vorangegangenen Untersuchungen}{section-six}
\begin{center}\Large SECHSTER ABSCHNITT.\par\medskip\normalsize VERSCHIEDENE ANWENDUNGEN DER VORANGEGANGENEN UNTERSUCHUNGEN.\end{center}\BookRule{0.16\linewidth}

\Needspace{7\baselineskip}
\Article{308}
Wie reich die höhere Arithmetik an Wahrheiten ist, die auch in anderen Teilen der Mathematik nützlich sind, haben wir bereits an mehreren Stellen angedeutet. Einige Anwendungen, die eine ausführlichere Darstellung verdienen, gesondert zu behandeln, schien uns nicht unnütz. Dabei wollen wir diesen Gegenstand, der leicht mehrere Bände füllen könnte, weniger erschöpfen als vielmehr durch einige Beispiele erläutern. In diesem Abschnitt behandeln wir zuerst die Zerlegung von Brüchen in einfachere, dann die Umwandlung gewöhnlicher Brüche in Dezimalbrüche. Danach erklären wir ein neues Ausschlussverfahren zur Lösung unbestimmter Gleichungen zweiten Grades. Schließlich geben wir neue, bequeme Methoden an, Primzahlen von zusammengesetzten Zahlen zu unterscheiden und die Faktoren der letzten zu ermitteln. Im folgenden Abschnitt begründen wir die allgemeine Theorie einer besonderen Art von Funktionen, die sich durch die ganze Analysis sehr weit erstreckt, soweit sie mit der höheren Arithmetik sehr eng zusammenhängt. Insbesondere wollen wir die Theorie der Kreisteilung, von der bisher nur die ersten Grundlagen bekannt geworden sind, durch neue Ergebnisse erweitern.

\Needspace{9\baselineskip}\begin{center}\textit{Zerlegung von Brüchen in einfachere.}\end{center}

\Needspace{7\baselineskip}
\Article{309}
\textsc{Aufgabe.} \textit{Einen Bruch \(\frac mn\), dessen Nenner \(n\) das Produkt zweier teilerfremder Zahlen \(a,b\) ist, in zwei andere mit den Nennern \(a,b\) zerlegen.}

\SourceLeaf{386}
\textit{Lösung.} Die gesuchten Brüche seien \(\frac xa,\frac yb\). Dann muss \(bx+ay=m\) gelten. Also ist \(x\) eine Wurzel der Kongruenz \(bx\equiv m\; (\mathrm{mod}.\ a)\), die sich nach Abschnitt II finden lässt, während dann \(y\)\(=\frac{m-bx}{a}\) gilt.

Bekanntlich hat die Kongruenz \(bx\equiv m\) unendlich viele Wurzeln, die aber nach \(a\) kongruent sind, und nur eine positive, die kleiner als \(a\) ist. Es kann allerdings auch vorkommen, dass \(y\) negativ wird. Wir brauchen kaum zu bemerken, dass \(y\) auch durch die Kongruenz \(ay\equiv m\; (\mathrm{mod}.\ b)\) und \(x\) durch die Gleichung \(x=\frac{m-ay}{b}\) gefunden werden kann. — \textit{Zum Beispiel} ist für den gegebenen Bruch \(\frac{58}{77}\) die Zahl \(4\) der Wert des Ausdrucks \(\frac{58}{11}\; (\mathrm{mod}.\ 7)\). Daher zerlegt sich \(\frac{58}{77}\) in \(\frac47+\frac2{11}\).

\Needspace{7\baselineskip}
\Article{310}
Ist ein Bruch \(\frac mn\) gegeben, dessen Nenner \(n\) das Produkt beliebig vieler paarweise teilerfremder Faktoren \(a,b,c,d\) usw. ist, so kann er nach dem vorigen Artikel zunächst in zwei mit den Nennern \(a\) und \(bcd\) usw. zerlegt werden; der zweite wieder in zwei mit den Nennern \(b\) und \(cd\) usw.; der letzte wiederum in zwei und so fort. Schließlich erhält der gegebene Bruch die Form \[\frac mn=\frac\alpha a+\frac\beta b+\frac\gamma c+\frac\delta d+\text{usw.}\]
Die Zähler \(\alpha,\beta,\gamma,\delta\) usw. können offenbar positiv und kleiner als ihre jeweiligen Nenner gewählt werden, außer dem letzten. Sind die übrigen bestimmt, so ist der letzte nicht mehr beliebig und kann auch negativ oder größer als sein Nenner werden, da wir \(m<n\) nicht voraussetzen. Dann ist es meist zweckmäßig, ihn in die Form \(\frac\varepsilon e\mp k\) zu bringen, wobei \(\varepsilon\) positiv und kleiner als \(e\), dagegen \(k\) eine ganze Zahl ist. Schließlich können \(a,b,c\) usw. offenbar als Primzahlen oder Primzahlpotenzen gewählt werden.

\textit{Beispiel.} Der Bruch \(\frac{391}{924}\), mit dem Nenner \(=4.3.7.11\), zerlegt sich auf diese Weise in \(\frac14+\frac{40}{231}\); \(\frac{40}{231}\) in \(\frac23-\frac{38}{77}\); \(-\frac{38}{77}\) in \(\frac17-\frac7{11}\). Schreibt man daher \(\frac4{11}-1\) für \(-\frac7{11}\), so erhält man \(\frac{391}{924}=\frac14+\frac23+\frac17+\frac4{11}-1\).

\Needspace{7\baselineskip}
\Article{311}
Ein Bruch \(\frac mn\) lässt sich \textit{nur auf eine Weise} in die Form \(\frac\alpha a+\frac\beta b+\text{usw.}\mp k\) bringen, sodass \(\alpha,\beta\) usw. positiv und kleiner als \(a,b\) usw. sind. Setzen wir nämlich \[\frac mn=\frac\alpha a+\frac\beta b+\frac\gamma c+\text{usw.}\mp k=\frac{\alpha'}a+\frac{\beta'}b+\frac{\gamma'}c+\text{usw.}\mp k'\]
\SourceLeaf{387}
und auch \(\alpha',\beta'\) usw. positiv und kleiner als \(a,b\) usw. voraus, so gilt notwendig \(\alpha=\alpha',\ \beta=\beta',\ \gamma=\gamma'\) usw. und \(k=k'\). Denn durch Multiplikation mit \(n=abc\) usw. ergibt sich offenbar \(m\equiv\alpha bcd\text{ usw.}\equiv\alpha'bcd\text{ usw.}\; (\mathrm{mod}.\ a)\). Da \(bcd\) usw. zu \(a\) teilerfremd ist, gilt notwendig \(\alpha\equiv\alpha'\), also \(\alpha=\alpha'\), und ebenso \(\beta=\beta'\) usw. Daraus folgt auch sofort \(k=k'\). Da nun völlig beliebig ist, zu welchem Nenner der Zähler zuerst berechnet wird, lassen sich offenbar \textit{alle} Zähler ebenso bestimmen wie \(\alpha\) im vorigen Artikel, etwa \(\beta\) durch die Kongruenz \(\beta acd\text{ usw.}\equiv m\; (\mathrm{mod}.\ b)\), \(\gamma\) durch \(\gamma abd\text{ usw.}\equiv m\; (\mathrm{mod}.\ c)\) usw. Die Summe aller so gefundenen Brüche ist entweder gleich dem gegebenen Bruch \(\frac mn\), oder die Differenz ist eine ganze Zahl \(=k\). Auf diese Weise erhalten wir zugleich eine Kontrolle der Rechnung. So liefern im Beispiel des vorigen Artikels die Werte der Ausdrücke \(\frac{391}{231}\; (\mathrm{mod}.\ 4)\), \(\frac{391}{308}\; (\mathrm{mod}.\ 3)\), \(\frac{391}{132}\; (\mathrm{mod}.\ 7)\), \(\frac{391}{84}\; (\mathrm{mod}.\ 11)\) unmittelbar die Zähler \(1,2,1,4\) zu den Nennern \(4,3,7,11\). Die Summe dieser Brüche übersteigt, wie man findet, den gegebenen Bruch um eins.
